\chapter{RAG Corporativo e Arquitetura de Dados}
\label{chap:rag_corporativo}

% ---------------------------------------------------------------
\section{Escopo do capitulo}
% ---------------------------------------------------------------

Este capitulo apresenta, em formato de aula-laboratorio, a implementacao de
um pipeline \textit{Retrieval-Augmented Generation} (RAG\index{RAG|textbf})
para contexto bancario, com foco em fontes heterogeneas,
indexacao\index{indexacao (RAG)|textbf}, recuperacao\index{recuperacao semantica}
e governanca de dados. O objetivo pratico consiste em construir um fluxo no
qual documentos corporativos e entrevistas de clientes sao convertidos em
contexto recuperavel para agentes, reduzindo alucinacao e aumentando
fidelidade da resposta.

% ---------------------------------------------------------------
\section{Objetivos de aprendizagem}
% ---------------------------------------------------------------

Ao concluir este capitulo, o leitor sera capaz de:

\begin{itemize}[leftmargin=*]
  \item Diferenciar conhecimento parametricamente aprendido pelo LLM e
        conhecimento recuperado por RAG;
  \item Modelar documentos de negocio para indexacao em
        OpenSearch\index{OpenSearch|textbf};
  \item Implementar ingestao textual para os tipos
        \texttt{documento} e \texttt{entrevista};
  \item Validar, com casos de teste, se a resposta do agente utiliza o
        contexto recuperado de forma verificavel.
\end{itemize}

% ---------------------------------------------------------------
\section{Conceitos teoricos}
% ---------------------------------------------------------------

\subsection{RAG em ambiente regulado}

Em ambiente bancario, perguntas de negocio dependem de politicas, regras de
produto, condicoes de canal e historico de relacionamento. Essas informacoes
mudam com frequencia e nao devem depender exclusivamente do conhecimento
interno do modelo. O padrao RAG injeta contexto em tempo de execucao,
desacoplando atualizacao documental de novo \textit{fine-tuning}.

\subsection{Tipos de documento no projeto}

O projeto adota dois tipos principais:

\begin{itemize}[leftmargin=*]
  \item \textbf{\texttt{documento}:} politicas, FAQ, regras de oferta e
        materiais normativos;
  \item \textbf{\texttt{entrevista}:} texto de voz do cliente ou de grupo,
        utilizado para os modos \textit{persona} e \textit{twin}.
\end{itemize}

Essa separacao permite governanca por origem e facilita auditoria de qual
fonte sustentou determinada resposta.

% ---------------------------------------------------------------
\section{Implementacao orientada por passos}
% ---------------------------------------------------------------

\subsection{Passo 1 --- Preparar colecao de documentos}

Nesta etapa, organiza-se um conjunto minimo para reproducao da aula:

\begin{lstlisting}[language=bash,
  caption={Colecao minima para laboratorio de RAG corporativo.},
  label={lst:rag_dataset_minimo}]
dados_rag/
  documentos/
    politica_investimentos_2026.txt
    faq_cartao_credito.txt
  entrevistas/
    cliente_C00042.txt
    cluster_1_focus_group.txt
\end{lstlisting}

Cada arquivo deve conter texto continuo, sem metadados sensiveis. Campos de
identificacao (cliente ou cluster) devem ser enviados explicitamente no payload.

\subsection{Passo 2 --- Ingerir texto no endpoint \texttt{/documentos}}

O endpoint de ingestao esta definido no contrato OpenAPI e aceita JSON.
No laboratorio, utiliza-se a chamada abaixo:

\begin{lstlisting}[language=bash,
  caption={Exemplo de ingestao de documento de politica no RAG.},
  label={lst:ingestao_documento_api}]
curl -X POST "$API_URL/documentos" \
  -H "Content-Type: application/json" \
  -d '{
    "tipo": "documento",
    "titulo": "Politica de Investimentos 2026",
    "texto": "Clientes conservadores priorizam liquidez diaria..."
  }'
\end{lstlisting}

Para entrevista individual:

\begin{lstlisting}[language=bash,
  caption={Ingestao de entrevista associada a cliente.},
  label={lst:ingestao_entrevista_cliente}]
curl -X POST "$API_URL/documentos" \
  -H "Content-Type: application/json" \
  -d '{
    "tipo": "entrevista",
    "titulo": "Entrevista cliente C00042",
    "cliente_id": "C00042",
    "texto": "Prefiro resolver tudo no app e evitar taxas altas..."
  }'
\end{lstlisting}

\subsection{Passo 3 --- Verificar chunking e indexacao}

A funcao de ingestao fragmenta textos em \textit{chunks}\index{chunking|textbf}
e os grava em indices OpenSearch distintos para cada finalidade.
Conceitualmente, o fluxo corresponde a:

\begin{lstlisting}[language=Python,
  caption={Fluxo simplificado da indexacao no ingester.},
  label={lst:ingester_fluxo_simplificado}]
if tipo == "entrevista":
    chunks = chunkar(texto)
    indexar("entrevistas-clientes", chunks, metadata)
else:
    chunks = chunkar(texto)
    indexar("documentos-knowledge", chunks, metadata)
\end{lstlisting}

O resultado esperado da API e um JSON contendo \texttt{chunks\_indexados > 0}.

\subsection{Passo 4 --- Consumir o contexto em consulta real}

Apos indexar, executa-se uma pergunta de negocio no endpoint ass\'incrono:

\begin{lstlisting}[language=bash,
  caption={Consulta de negocio para verificar uso do contexto RAG.},
  label={lst:consulta_rag_laboratorio}]
curl -X POST "$API_URL/query" \
  -H "Content-Type: application/json" \
  -d '{
    "modo": "segmento",
    "pergunta": "Qual produto devo oferecer para perfil conservador?",
    "dados_cliente": {
      "idade": 55,
      "renda_mensal": 18000,
      "saldo_medio": 120000,
      "transacoes_mes": 25,
      "score_credito": 820,
      "num_produtos": 7
    }
  }'
\end{lstlisting}

Em seguida, consulta-se \texttt{/status/\{request\_id\}} ate \texttt{COMPLETED}
e compara-se a resposta com os documentos enviados.

\subsection{Passo 5 --- Resultado completo da aula}

Ao final da pratica, o ambiente deve conter:

\begin{itemize}[leftmargin=*]
  \item indice \texttt{documentos-knowledge} com politicas e FAQ;
  \item indice \texttt{entrevistas-clientes} com voz de cliente/grupo;
  \item resposta do agente sustentada por contexto recuperado;
  \item evidencias de payload, \texttt{request\_id} e resposta final arquivadas.
\end{itemize}

% ---------------------------------------------------------------
\section{Plano de testes do capitulo}
% ---------------------------------------------------------------

\subsection{Teste funcional}

\begin{itemize}[leftmargin=*]
  \item \textbf{TC-01:} Ingerir um \texttt{documento} e verificar
        \texttt{statusCode = 200} e \texttt{chunks\_indexados > 0}.
  \item \textbf{TC-02:} Ingerir uma \texttt{entrevista} com
        \texttt{cliente\_id} e validar persistencia no indice de entrevistas.
  \item \textbf{TC-03:} Executar consulta \texttt{modo=segmento} e confirmar
        resposta final sem erro de processamento.
\end{itemize}

\subsection{Teste de desempenho}

\begin{itemize}[leftmargin=*]
  \item \textbf{TP-01:} Medir tempo de ingestao de 10 documentos de 2 KB;
        criterio: mediana $\leq 2{,}5$\,s por documento.
  \item \textbf{TP-02:} Medir latencia de resposta ass\'incrona ponta a ponta;
        criterio: p95 $\leq 20$\,s em ambiente \texttt{dev}.
\end{itemize}

\subsection{Teste de qualidade de resposta}

\begin{itemize}[leftmargin=*]
  \item \textbf{TQ-01:} Verificar, em 10 perguntas, se a resposta cita
        explicitamente contexto compativel com os documentos indexados.
  \item \textbf{TQ-02:} Avaliar fidelidade com LLM-as-Judge (Capitulo~\ref{chap:llm_as_judge});
        criterio: media $\geq 8{,}0/10$.
\end{itemize}

% ---------------------------------------------------------------
\section{Criterios de aceite}
% ---------------------------------------------------------------

\begin{itemize}[leftmargin=*]
  \item TC-01 a TC-03 aprovados sem erro bloqueante.
  \item TP-01 e TP-02 dentro dos limiares definidos.
  \item TQ-01 com aderencia documental observavel em pelo menos 80\% dos casos.
\end{itemize}

% ---------------------------------------------------------------
\section{Espaco para expansao iterativa}
% ---------------------------------------------------------------

\begin{itemize}[leftmargin=*]
  \item \textbf{v6.1} --- Adicionar suporte a PDF/DOCX com pipeline de extracao.
  \item \textbf{v6.2} --- Incluir versionamento de documento e expiracao por vigencia.
  \item \textbf{Lacuna} --- Investigar \textit{reranking}\index{reranking} para aumentar precisao top-k.
\end{itemize}
